O programa Apollo (1961--1972) consolidou capacidades operacionais diretamente relevantes para \gls{rvdb}, em especial na condução de fases de aproximação e acoplamento que exigem navegação orbital relativa, guiamento e tomada de decisão em tempo real \cite{RVDBapollo}.
Nesse contexto, a interação homem--máquina evoluiu de comandos diretos para um papel de supervisão, principalmente no monitoramento de estados, seleção de modos e gerenciamento de contingências, o que reduziu a dependência contínua da estação terrestre \cite{nevins1971}.

Como descrito por \cite{nevins1971}, requisitos associados a essa autonomia operacional incluíam:
\begin{enumerate}[label=(\roman*)]
    \item o sistema deveria ser capaz de completar a missão sem auxílio da estação terrestre;
    \item o sistema empregaria a participação humana sempre que isso pudesse simplificar ou melhorar a operação em comparação às sequências pré-estabelecidas das funções exigidas;
    \item o sistema deveria fornecer interfaces adequadas ao piloto por meio das telas e dos métodos de controle do sistema de orientação;
    \item o sistema deveria ser projetado de modo que as tarefas pudessem ser executadas por apenas um operador, de forma segura, permitindo o retorno seguro à Terra a partir de qualquer ponto da missão.
\end{enumerate}

Em operações de \gls{rvdb}, esses requisitos se traduzem na necessidade de (i) interpretar informações de navegação e orientação, (ii) intervir quando necessário em decisões de guiamento e controle e (iii) garantir continuidade operacional por meio de modos redundantes e transições seguras entre sistemas primários e de \emph{backup} \cite{nevins1971}.
